\appendix

\section{Proof Details}

\subsection{Coding localization, grid rounding, and recentering}

\begin{proposition}[Verified non-adaptive localization and scalar wrapper]
\label{prop:localization}
Under Assumptions~\ref{assump:parameter-domain}--
\ref{assump:precommitted-protocol}, set $\eta=\delta/4$.  There is a
deterministic precommitted localization block using one Borel-set bit per
sample and a decoder defined on every transcript such that, with
$c=\operatorname{mid}(I)$ and $\mu=\mu(D)$,
\[
 \mathcal E_{\rm loc}:=\{|c-\mu|\leq50\sigma\},
 \qquad \Pr_D(\mathcal E_{\rm loc})\geq1-\frac\delta4.
\]
Moreover, with $C_{{\rm loc},k}=10001$,
\[
 N_{\rm loc}\leq C_{{\rm loc},k}
 \left[1+\log\frac\lambda\sigma+\log\frac4\delta\right].
\]
In the trivial source branch $I=[-\lambda,\lambda]$, $c=0$, and
$N_{\rm loc}=0$.  In the nontrivial branch the fixed minimum-index Hamming
tie rule, $I$, and $c$ remain defined even when localization fails.
\end{proposition}

\begin{proof}
We instantiate Theorem~16 of \citet{lau-scarlett-2026}, whose source label is
\texttt{thm: alternative localization}, together with the construction in
the source appendix labeled \texttt{appendix: two-stage}.  That result assumes
$\lambda\geq\sigma>0$, $\eta\in(0,1/2)$, iid observations with
$\mu\in[-\lambda,\lambda]$, and $\mathbb E|X-\mu|\leq\sigma$.
Assumptions~\ref{assump:parameter-domain} and \ref{assump:moment-class}
give all but the last condition.  Holder's inequality supplies it directly:
\[
 \mathbb E|X-\mu|
 \leq\bigl(\mathbb E|X-\mu|^k\bigr)^{1/k}
       \bigl(\mathbb E1^{k/(k-1)}\bigr)^{(k-1)/k}
 \leq\sigma.
\]
Assumption~\ref{assump:independent-samples} gives the iid source block, and
Assumption~\ref{assump:precommitted-protocol} gives its fixed placement in
the full protocol.

For completeness, the numerical source interface is as follows.  Put
$h_{\rm src}=20\sigma$.  If $2\lambda\leq h_{\rm src}$, the construction
asks no query and returns $I=[-\lambda,\lambda]$.  Thus
\[
 |I|=2\lambda\leq20\sigma,\qquad
 |c-\mu|=|\mu|\leq\lambda\leq10\sigma,
\]
so $\mathcal E_{\rm loc}$ occurs deterministically and $N_{\rm loc}=0$.

If $2\lambda>h_{\rm src}$, define
\[
 N_{\rm src}=\left\lceil\frac{2\lambda}{h_{\rm src}}\right\rceil,
 \qquad
 \Delta=\frac{2\lambda}{N_{\rm src}},
 \qquad
 \ell=\left\lceil10000\left(
    \log N_{\rm src}+\log\frac1\eta\right)\right\rceil.
\]
Then $h_{\rm src}/2\leq\Delta\leq h_{\rm src}$.  With
$e_j^{\rm src}=-\lambda+(j-1)\Delta$, use the half-open bins
$K_j^{\rm src}=[e_j^{\rm src},e_{j+1}^{\rm src})$, closing only the final
right endpoint, and clip observations outside $[-\lambda,\lambda]$ to the
first or last bin.  Hence every boundary has a unique bin index.  The source
fixes a deterministic balanced binary codeword of length $\ell$ for every
bin, queries coordinate $t$ of the codeword of the clipped bin containing
$X_t$, decodes a minimum-Hamming-score bin with the minimum-index tie rule,
and returns the union of that bin and at most its two neighbors on each side.
The interval is therefore defined on every word and obeys
\[
 |I|\leq5\Delta\leq100\sigma.
\]
The source guarantee gives
$\Pr_D\{\mu\in I\}\geq1-\eta=1-\delta/4$.  On $\{\mu\in I\}$,
\[
 |c-\mu|\leq\frac{|I|}{2}\leq50\sigma,
\]
which proves the required probability statement.  The source's
location-dependent refinement stage is not used.

To map each source bit into our query model, fix the promised codebook, take
$R_{\rm loc}$ to be degenerate, and put
$\mathcal B_t=Q_t^{-1}(\{1\})$.  The clipped-bin map is constant on finitely
many half-open or closed intervals and on the two outer rays; hence every
$\mathcal B_t$ is Borel and
$Y_t^{\rm loc}=\mathbf1\{X_t\in\mathcal B_t\}=Q_t(X_t)$.  The bins,
queries, codebook, sample count, and tie rule depend only on the known
parameters and are fixed before any response.

Finally, in the nontrivial branch
\[
 N_{\rm src}=\left\lceil\frac\lambda{10\sigma}\right\rceil\geq2.
\]
For $x=\lambda/(10\sigma)>1$,
\[
 N_{\rm src}=\lceil x\rceil\leq x+1\leq2x
 =\frac\lambda{5\sigma}\leq\frac\lambda\sigma,
\]
and therefore
\[
\begin{aligned}
 N_{\rm loc}
 &\leq1+10000\left(\log\frac\lambda\sigma+log\frac4\delta\right)\\
 &\leq10001\left(1+\log\frac\lambda\sigma+log\frac4\delta\right).
\end{aligned}
\]
This proves the proposition in both branches, including every tie and failure
transcript.
\end{proof}

\begin{lemma}[Midpoint-to-grid core bridge]\label{lem:localization-core}
Under Assumption~\ref{assump:parameter-domain} and
Proposition~\ref{prop:localization}, define $h_0=a_k\sigma$ and use the
minimum-index nearest-center rule.  Every realized $c$ satisfies
\[
 |c-m_0|\leq\frac{h_0}{2}.
\]
If also $|c-\mu(D)|\leq50\sigma$ and $a_k\geq200$, then
\[
 |m_0-\mu(D)|\leq\frac{3h_0}{4},\qquad
 \mu(D)\in[m_0-3h_0/4,m_0+3h_0/4]
 \subseteq J_{0,j_0(c)}.
\]
The conclusion includes both source branches and every grid tie.
\end{lemma}

\begin{proof}
Let $q=\lfloor c/h_0\rfloor$.  The candidate
$m_{0,q}=(q+1/2)h_0$ is at distance at most $h_0/2$ from $c$, so the same
is true of a nearest center.  At a Voronoi boundary both candidates are at
exactly that distance and the smaller-index rule preserves the bound.  In
fact
\[
 j_0(c)=\left\lceil\frac c{h_0}\right\rceil-1,
\]
so the rule is Borel and also covers negative $c$.

On $\mathcal E_{\rm loc}$, Proposition~\ref{prop:localization} and
$a_k\geq200$ give
\[
\begin{aligned}
 |m_0-\mu(D)|
 &\leq|m_0-c|+|c-\mu(D)|\\
 &\leq\frac{h_0}{2}+50\sigma
 =\left(\frac12+\frac{50}{a_k}\right)h_0
 \leq\frac{3h_0}{4}.
\end{aligned}
\]
The closed core is contained in
$J_{0,j_0}=[m_0-3h_0/2,m_0+3h_0/2)$, including its two core endpoints.
\end{proof}

\begin{lemma}[Recentered moment certificate]\label{lem:recentered-moment}
Under Assumptions~\ref{assump:parameter-domain} and
\ref{assump:moment-class} and Lemma~\ref{lem:localization-core}, if
$|c-\mu(D)|\leq50\sigma$ and $a_k\geq200$, then
\[
 M_k(c):=\int_{\mathbb R}|x-m_0(c)|^kD(dx)
 \leq C_k^{\rm rec}\sigma^k,
 \qquad
 C_k^{\rm rec}:=2^{k-1}\left[1+\left(\frac{3a_k}{4}\right)^k\right].
\]
In particular, the recentered moment is finite and its constant depends only
on the fixed $k$.
\end{lemma}

\begin{proof}
For fixed $c$ on $\mathcal E_{\rm loc}$, the power-triangle inequality gives
\[
 |x-m_0|^k
 \leq2^{k-1}\bigl(|x-\mu(D)|^k+|\mu(D)-m_0|^k\bigr).
\]
Indeed, this follows from convexity because
$(|u|+|v|)^k\leq2^{k-1}(|u|^k+|v|^k)$.  Integrating and applying
Assumption~\ref{assump:moment-class} and Lemma~
\ref{lem:localization-core} yields
\[
\begin{aligned}
 M_k(c)
 &\leq2^{k-1}\left(\mathbb E_D|X-\mu(D)|^k
          +|\mu(D)-m_0|^k\right)\\
 &\leq2^{k-1}\left[\sigma^k+(3h_0/4)^k\right]
 =C_k^{\rm rec}\sigma^k.
\end{aligned}
\]
\end{proof}

\begin{proposition}[Independent refinement product law]
\label{prop:refinement-independence}
Under Assumptions~\ref{assump:parameter-domain}--
\ref{assump:precommitted-protocol} and Proposition~\ref{prop:localization},
let $\mathscr L_{\rm loc}$ be the sigma-field generated by the complete
localization block, including its local samples, seed, fixed queries, bits,
interval $I$, and outputs $c,m_0$.  It is independent of the sigma-field
generated by all refinement samples, levels, colors, branches, types, masks,
dithers, and group assignments.  The variables $c,m_0$,
$\mathcal E_{\rm loc}$ and, for fixed $D$, $M_k(c)$ are
$\mathscr L_{\rm loc}$-measurable.  Conditional on
$\mathscr L_{\rm loc}$, the per-sample refinement tuples retain their
original independent product law; for every bounded measurable function $g$
of the complete refinement block,
\[
 \mathbb E_D[g\mid\mathscr L_{\rm loc}]=\mathbb E_Dg
 \quad\text{almost surely}.
\]
\end{proposition}

\begin{proof}
Every localization query, bit, $I$, $c$, and $m_0$ is a measurable function
of the localization samples and the deterministic source objects.  The event
$\mathcal E_{\rm loc}$ is therefore $\mathscr L_{\rm loc}$-measurable.
For fixed $D$, consider
\[
 \Phi_D(m)=\left(\int_{\mathbb R}|x-m|^kD(dx)\right)^{1/k}.
\]
Minkowski's inequality and Assumption~\ref{assump:moment-class} give
\[
 \Phi_D(m)\leq\|X-\mu(D)\|_{L_k(D)}+|m-\mu(D)|
 \leq\sigma+|m-\mu(D)|<\infty.
\]
The reverse triangle inequality in $L_k(D)$ gives, for all $m,m'\in\mathbb R$,
\[
 |\Phi_D(m)-\Phi_D(m')|
 \leq\|(X-m)-(X-m')\|_{L_k(D)}=|m-m'|.
\]
Thus $\Phi_D$ is Borel.  Since $c\mapsto m_0(c)$ is Borel by
Lemma~\ref{lem:localization-core},
$M_k(c)=\Phi_D(m_0(c))^k$ is $\mathscr L_{\rm loc}$-measurable.

Assumption~\ref{assump:independent-samples} makes the complete refinement
sample-and-seed family independent of the complete localization block.
Measurable localization transformations preserve that independence.
Conditioning on the first sigma-field therefore leaves the second sigma-field's
joint law unchanged; this law is a product across sample indices by the same
assumption and the precommitted construction.  The conditional expectation
identity follows, including on localization-failure transcripts.

Together, Proposition~\ref{prop:localization}, Lemmas~
\ref{lem:localization-core} and \ref{lem:recentered-moment}, and this product
law give the complete localization interface: midpoint radius $50\sigma$,
the selected-core certificate, the explicit recentered moment, and an
independent refinement array, with the source and grid ceilings already
included in $N_{\rm loc}$.
\end{proof}
\subsection{Dyadic ancestry, ring partition, and Borel query geometry}

\begin{lemma}[Dyadic ancestor and outer-radius geometry]
\label{lem:dyadic-ancestor}
Under Assumption~\ref{assump:parameter-domain} and
Proposition~\ref{prop:localization}, let $c\in\mathbb R$ be any realized
decoder value and define $j_0,j_s,b_s,m_s,d_s$ as in
Section~\ref{sec:preliminaries}.  For every $1\leq s\leq S$,
\[
 j_{s-1}=2j_s+b_s,qquad b_s\in\{0,1\},qquad
 J_{s-1,j_{s-1}}\subseteq J_{s,j_s}.
\]
Moreover,
\[
 |d_s|<\frac{h_s}{2}\leq h_s,qquad
 [m_0-H,m_0+H]\subset J_{S,j_S}.
\]
These statements hold for positive and negative indices, both child values,
and every nearest-center or half-open-grid tie.
\end{lemma}

\begin{proof}
Put $t=c/h_0$.  The integer $j_0=\lceil t\rceil-1$ satisfies
$j_0<t\leq j_0+1$.  If the last inequality is strict, its center is the
unique nearest center; if equality holds, the centers indexed by $j_0$ and
$j_0+1$ tie and $j_0$ is selected.  This covers negative $t$ as well.

For $s\geq1$ define the Euclidean remainder
\[
 r_s=j_0-2^s\left\lfloor\frac{j_0}{2^s}\right\rfloor
     =j_0-2^sj_s.
\]
The floor inequalities give $r_s\in\{0,1,\ldots,2^s-1\}$, including for
$j_0<0$, and hence
\[
\begin{aligned}
 j_{s-1}
 &=\left\lfloor\frac{2^sj_s+r_s}{2^{s-1}}\right\rfloor
 =2j_s+\left\lfloor\frac{r_s}{2^{s-1}}\right\rfloor,\\
 b_s&=\left\lfloor\frac{r_s}{2^{s-1}}\right\rfloor\in\{0,1\}.
\end{aligned}
\]
Since $h_s=2h_{s-1}$, direct endpoint calculation gives
\[
\begin{aligned}
 J_{s,j}&=[(j-1)h_s,(j+2)h_s),\\
 J_{s-1,2j}&=[(j-\tfrac12)h_s,(j+1)h_s),\\
 J_{s-1,2j+1}&=[jh_s,(j+\tfrac32)h_s).
\end{aligned}
\]
Both child paddings lie in $J_{s,j}$.  Substitution of $j_s,b_s$ and
iteration prove
$J_{0,j_0}\subseteq\cdots\subseteq J_{S,j_S}$.

The same remainder identity yields
\[
 m_0=(j_s+\vartheta_s)h_s,qquad
 \vartheta_s=\frac{r_s+1/2}{2^s}\in(0,1).
\]
Consequently
\[
 d_s=(\tfrac12-\vartheta_s)h_s,qquad |d_s|<h_s/2.
\]
At $s=S$, $H=h_S$ and
\[
 (j_S-1)H<m_0-H\leq m_0+H<(j_S+2)H.
\]
Thus both endpoints of the closed radius-$H$ interval lie in the half-open
padding $J_{S,j_S}$.
\end{proof}

\begin{proposition}[Exact padded-ring partition]\label{prop:ring-partition}
Under Assumption~\ref{assump:parameter-domain} and
Lemma~\ref{lem:dyadic-ancestor}, for every $c\in\mathbb R$ the target rings
are pairwise disjoint and
\[
 \mathop{\dot\bigcup}_{s=0}^S R_s(c)=J_{S,j_S},
 \qquad
 \sum_{s=0}^S\mathbf1\{x\in R_s(c)\}
 =\mathbf1\{x\in J_{S,j_S}\}
\]
for every $x\in\mathbb R$.  Every $x\in J_{0,j_0}$ belongs to no target
ring $R_s(c)$ with $s\geq1$.
\end{proposition}

\begin{proof}
Write $\mathbf J_s=J_{s,j_s}$.  The exact child identity in
Lemma~\ref{lem:dyadic-ancestor} gives, for $s\geq1$,
\[
 R_s(c)=\mathcal R_{s,j_s,b_s}
 =J_{s,j_s}\setminus J_{s-1,2j_s+b_s}
 =\mathbf J_s\setminus\mathbf J_{s-1},
\]
whereas $R_0(c)=\mathbf J_0$.  Therefore
\[
 \mathbf J_s=\mathbf J_{s-1}\mathbin{\dot\cup}R_s(c).
\]
Iterating this identity yields the asserted disjoint union.  Equivalently,
if $t<s$, then
$R_t(c)\subseteq\mathbf J_t\subseteq\mathbf J_{s-1}$, whereas
$R_s(c)\cap\mathbf J_{s-1}=\varnothing$.  This also proves the pointwise
indicator identity at every half-open endpoint.  Finally, if
$x\in\mathbf J_0$, nesting puts $x$ in every $\mathbf J_{s-1}$, so it is
in none of the differences $\mathbf J_s\setminus\mathbf J_{s-1}$.
\end{proof}

\begin{lemma}[Four-color support, endpoint, and amplitude geometry]
\label{lem:four-color}
Under Assumption~\ref{assump:parameter-domain} and
Lemma~\ref{lem:dyadic-ancestor}, fix $0\leq s\leq S$ and
$\ell\in\{0,1,2,3\}$.  The paddings
$\{J_{s,j}:j\equiv\ell\pmod4\}$ are pairwise disjoint, and so are the
rings at any fixed branch.  All these sets are Borel and
\[
 \left|\frac{x-m_{s,j}}{2h_s}\right|\leq\frac34
 \quad\text{on every such ring},
\]
while the mass amplitude is one.  If $s\geq1$ and
$x\in J_{0,j_0}$, then $x$ lies in neither the target ring nor any alias
$\mathcal R_{s,j,b_s}$ with $j\equiv j_s\pmod4$.
\end{lemma}

\begin{proof}
The padding is $J_{s,j}=[(j-1)h_s,(j+2)h_s)$.  If $j<j'$ have the same
residue modulo four, then $j'\geq j+4$, so the right endpoint of the first
padding is $(j+2)h_s$ and the left endpoint of the second is at least
$(j+3)h_s$.  They are disjoint with a gap of at least one cell.  This
calculation is unchanged for negative indices.

For $s\geq1$, the child-padding formulas give
\[
\begin{aligned}
 \mathcal R_{s,j,0}
 &=[(j-1)h_s,(j-\tfrac12)h_s)
   \cup[(j+1)h_s,(j+2)h_s),\\
 \mathcal R_{s,j,1}
 &=[(j-1)h_s,jh_s)
   \cup[(j+\tfrac32)h_s,(j+2)h_s).
\end{aligned}
\]
At level zero, $\mathcal R_{0,j}=J_{0,j}$.  Thus all sets are Borel and
the formulas settle each boundary.  For $x\in J_{s,j}$,
\[
 -\frac{3h_s}{2}\leq x-m_{s,j}<\frac{3h_s}{2},
\]
which proves the coordinate bound after division by $2h_s$; equality at the
included left endpoint is allowed.

For the final assertion, take $x\in J_{0,j_0}$.  Nesting puts it in
$J_{s,j_s}$ and, more strongly, in $J_{s-1,j_{s-1}}$, excluding it from
$\mathcal R_{s,j_s,b_s}$.  Every other $j\equiv j_s\pmod4$ has padding
disjoint from $J_{s,j_s}$, so its ring cannot contain $x$.  These cases
exhaust the retained color and branch.
\end{proof}

\begin{lemma}[Jointly Borel bounded precommitted query map]
\label{lem:query-measurability}
Under Assumptions~\ref{assump:parameter-domain} and
\ref{assump:precommitted-protocol} and Lemma~\ref{lem:four-color}, the
formula for each $F_i$ is jointly Borel in the sample and its complete
precommitted level, color, type, branch, and countable mask seed.  For every
seed realization and every $x\in\mathbb R$,
\[
 |F_i(x)|\leq1.
\]
Consequently each realized $F_i$ is Borel and every dithered superlevel set
$\{x:F_i(x)\geq u\}$, $u\in[-1,1]$, is Borel.  Neither map contains the
decoder output $c$.
\end{lemma}

\begin{proof}
For fixed $i$, let
\[
 \mathcal M_i=\{-1,+1\}^{\{0,\ldots,S\}\times\mathbb Z}
\]
carry its countable product Borel sigma-field.  The query-seed space is the
finite disjoint union
\[
\begin{aligned}
 \Omega_i^{\rm qry}
 :=&\ \bigl(\{0\}\times\{0,1,2,3\}\times\mathcal M_i\bigr)\\
 &\mathbin{\dot\cup}\mathop{\dot\bigcup}_{s=1}^S
 \bigl(\{s\}\times\{0,1,2,3\}
 \times\{{\rm coord},{\rm mass}\}\times\{0,1\}\times\mathcal M_i\bigr),
\end{aligned}
\]
with finite coordinates discrete.  The level-zero component correctly has no
child variable.

Fix $s,\ell$ and truncate the relevant sum to indices
$j\equiv\ell\pmod4$ with $|j|\leq N$.  Each summand is the product of a
mask coordinate, a continuous affine function or the constant one, and the
indicator of a Borel ring from Lemma~\ref{lem:four-color}.  Every finite
partial sum is jointly Borel.  Same-color disjointness implies that at each
$(x,\rho)$ at most one summand of the full countable sum is nonzero; the
partial sums therefore stabilize pointwise.  Their limit is jointly Borel.
Finite case distinctions over $L_i,C_i,T_i,B_i$ preserve this property.

The same uniqueness proves the range bound.  With no active ring the value is
zero; on a coordinate ring the unique summand has magnitude at most $3/4$;
on a mass ring it is a Rademacher sign and has magnitude one.  This covers
every mask realization and half-open endpoint without an absolute-summability
assumption.

For fixed seed $\omega$, $x\mapsto F_i(x,\omega)$ is Borel, and more
strongly
\[
 \{(x,\omega,u):F_i(x,\omega)-u\geq0\}
\]
is Borel in
$\mathbb R\times\Omega_i^{\rm qry}\times[-1,1]$.  Assumption~
\ref{assump:precommitted-protocol} places $\omega,u$ before the first response,
and the formula contains no $c$ or response bit.

Lemma~\ref{lem:dyadic-ancestor}, Proposition~\ref{prop:ring-partition},
Lemma~\ref{lem:four-color}, and this measurability result jointly give the
negative-index-safe nested path, exact ring partition, retained-ring
inactivity on $J_{0,j_0}$, and a bounded Borel query for every countable-mask
realization.
\end{proof}

\subsection{Dither inversion, mask projection, and the mean telescope}

\begin{lemma}[Exact centered-uniform dither moments]\label{lem:dither-moments}
Under Assumption~\ref{assump:independent-samples}, let
$U\sim{\rm Unif}[-1,1]$ and fix $f\in[-1,1]$ independently of $U$.  Define
\[
 \Delta_f(U)=\mathbf1\{f\geq U\}-\mathbf1\{0\geq U\}.
\]
Then, including $f=-1,0,1$,
\[
 \mathbb E_U\Delta_f(U)=\frac f2,
 \qquad
 \mathbb E_U\Delta_f(U)^2=\frac{|f|}{2}.
\]
If a random $F\in[-1,1]$ is independent of $U$, both identities hold
conditionally on $F$.
\end{lemma}

\begin{proof}
For $0\leq f\leq1$,
$\Delta_f(u)=\mathbf1_{(0,f]}(u)$, whereas for $-1\leq f<0$,
$\Delta_f(u)=-\mathbf1_{(f,0]}(u)$.  At $f=0$ the two indicators agree
identically.  Uniform measure on $[-1,1]$ has density $1/2$, so integrating
the signed interval gives $f/2$, while squaring removes its sign and gives
$|f|/2$.  The threshold endpoints are correct and have zero uniform mass.
Conditioning on an independent random $F$ fixes its value and applies the
same pointwise calculation.
\end{proof}

\begin{lemma}[Fresh-mask projection and exact alias cancellation]
\label{lem:mask-projection}
Under Assumption~\ref{assump:independent-samples}, Lemmas~
\ref{lem:dither-moments}, \ref{lem:four-color}, and
\ref{lem:query-measurability}, fix a level, color class, and, at a higher
level, branch and type.  Let $(E_j)_{j\in\mathcal J}$ be the resulting
pairwise-disjoint rings, let $(\rho_j)_{j\in\mathcal J}$ be independent
Rademacher masks, and write
\[
 F_\rho(x)=\sum_{j\in\mathcal J}
 \rho_j a_j(x)\mathbf1_{E_j}(x)\in[-1,1].
\]
For any target $j_\star\in\mathcal J$ and independent uniform dither $U$,
\[
 \mathbb E_{\rho,U}\left[
 \rho_{j_\star}\Delta_{F_\rho(x)}(U)\right]
 =\frac12a_{j_\star}(x)\mathbf1_{E_{j_\star}}(x)
\]
for every $x$.  In particular, a uniquely active alias $j\neq j_\star$
contributes exactly zero.
\end{lemma}

\begin{proof}
Conditioning first on the complete mask family and applying
Lemma~\ref{lem:dither-moments} gives
\[
 \mathbb E_{\rho,U}
 [\rho_{j_\star}\Delta_{F_\rho(x)}(U)]
 =\frac12\mathbb E_\rho[\rho_{j_\star}F_\rho(x)].
\]
At fixed $x$, Lemma~\ref{lem:four-color} leaves either no active summand or
one active index $q$.  In the first case both sides vanish.  In the second,
the last expectation is
\[
 \frac12a_q(x)\mathbf1_{E_q}(x)
 \mathbb E[\rho_{j_\star}\rho_q]
 =\frac12a_q(x)\mathbf1_{E_q}(x)\mathbf1\{q=j_\star\},
\]
because a mask squares to one and distinct masks are independent and
centered.  No countable sum is exchanged with expectation.
\end{proof}

\begin{proposition}[Exact per-level importance inversion]
\label{prop:level-inversion}
Under Assumptions~\ref{assump:independent-samples} and
\ref{assump:precommitted-protocol}, Propositions~
\ref{prop:refinement-independence} and \ref{prop:ring-partition}, and
Lemmas~\ref{lem:dyadic-ancestor}, \ref{lem:four-color},
\ref{lem:query-measurability}, and \ref{lem:mask-projection}, fix a decoder
value $c$ independent of the refinement block.  For every $x\in\mathbb R$,
put
\[
 K_s(x,c)=\mathbb E[W_i(c)\mathbf1\{L_i=s\}\mid X_i=x,c].
\]
Then
\[
 K_0(x,c)=(x-m_0)\mathbf1_{R_0(c)}(x),
\]
and, for $1\leq s\leq S$,
\[
 K_s(x,c)=\bigl[(x-m_s)+d_s\bigr]\mathbf1_{R_s(c)}(x)
 =(x-m_0)\mathbf1_{R_s(c)}(x).
\]
Every same-color alias contributes zero.
\end{proposition}

\begin{proof}
Fix $c,x$.  Proposition~\ref{prop:refinement-independence} leaves every
level and seed law unchanged and
\[
 \Delta Y_i
 =\mathbf1\{F_i(x)\geq U_i\}-\mathbf1\{0\geq U_i\}
 =\Delta_{F_i(x)}(U_i).
\]
At level zero, $W_i(c)$ vanishes unless $C_i=\kappa_0$, an event of
probability $1/4$.  Conditional on it, Lemma~\ref{lem:mask-projection}, with
target $j_0$ and amplitude $(x-m_{0,j})/(2h_0)$, yields
\[
 \mathbb E_{\rho,U}[\rho_{i,0,j_0}\Delta Y_i
 \mid x,c,L_i=0,C_i=\kappa_0]
 =\frac{x-m_0}{4h_0}\mathbf1_{R_0(c)}(x).
\]
All aliases have already vanished.  Substitution gives the complete ledger
\[
\begin{aligned}
 K_0(x,c)
 &=p_0\cdot\frac14\cdot\frac{16h_0}{p_0}
   \cdot\frac{x-m_0}{4h_0}\mathbf1_{R_0(c)}(x)\\
 &=(x-m_0)\mathbf1_{R_0(c)}(x).
\end{aligned}
\]
Thus the coefficient inverts the level, color, dither, and coordinate
normalization; no branch or type exists at level zero.

Now fix $s\geq1$.  The statistic vanishes unless
$C_i=\kappa_s$ and $B_i=b_s$, whose probabilities are $1/4$ and $1/2$.
For coordinate type, Lemma~\ref{lem:mask-projection} gives
\[
 \mathbb E_{\rho,U}[\rho_{i,s,j_s}\Delta Y_i
 \mid x,c,L_i=s,C_i=\kappa_s,B_i=b_s,T_i={\rm coord}]
 =\frac{x-m_s}{4h_s}\mathbf1_{R_s(c)}(x).
\]
For mass type it gives
\[
 \mathbb E_{\rho,U}[\rho_{i,s,j_s}\Delta Y_i
 \mid x,c,L_i=s,C_i=\kappa_s,B_i=b_s,T_i={\rm mass}]
 =\frac12\mathbf1_{R_s(c)}(x).
\]
Both expressions vanish for an active alias.  Since each type has probability
$1/2$, the coordinate contribution is
\[
\begin{aligned}
 &p_s\cdot\frac14\cdot\frac12\cdot\frac12
 \cdot\frac{16}{p_s}\cdot4h_s
 \cdot\frac{x-m_s}{4h_s}\mathbf1_{R_s(c)}(x)\\
 &\hspace{8em}=(x-m_s)\mathbf1_{R_s(c)}(x),
\end{aligned}
\]
and the mass contribution is
\[
\begin{aligned}
 &p_s\cdot\frac14\cdot\frac12\cdot\frac12
 \cdot\frac{16}{p_s}\cdot2d_s
 \cdot\frac12\mathbf1_{R_s(c)}(x)\\
 &\hspace{8em}=d_s\mathbf1_{R_s(c)}(x).
\end{aligned}
\]
These displays invert every level, color, branch, type, dither, and coordinate
factor separately.  Adding the mutually exclusive types and using
$d_s=m_s-m_0$ proves the result, including $d_s=0$.
\end{proof}

\begin{proposition}[Exact padded-ring mean telescope and residual]
\label{prop:mean-telescope}
Under Assumptions~\ref{assump:independent-samples} and
\ref{assump:precommitted-protocol}, Propositions~
\ref{prop:refinement-independence}, \ref{prop:ring-partition}, and
\ref{prop:level-inversion}, and Lemma~\ref{lem:four-color}, every decoder
value $c$ independent of the refinement block satisfies, for all $x$,
\[
 \mathbb E[W_i(c)\mid X_i=x,c]
 =(x-m_0)\mathbf1\{x\in J_{S,j_S}\}.
\]
Writing $\theta(c)=\mathbb E[W_i(c)\mid c]$, one has
\[
 \theta(c)=\mathbb E_D[(X-m_0)\mathbf1\{X\in J_{S,j_S}\}\mid c]
\]
and
\[
 \mu(D)-m_0=\theta(c)+
 \mathbb E_D[(X-m_0)\mathbf1\{X\notin J_{S,j_S}\}\mid c].
\]
There is no localization, fine-scale, ring-gap, or alias residual.  If
$D(J_{0,j_0})=1$, every retained higher correction is seedwise zero and the
outer residual is zero.  If $D\{m_0\}=1$, then $W_i(c)=0$ for every seed
realization.
\end{proposition}

\begin{proof}
The level events form a finite partition.  Proposition~
\ref{prop:level-inversion} and the exact ring partition give pointwise
\[
\begin{aligned}
 \mathbb E[W_i(c)\mid X_i=x,c]
 &=\sum_{s=0}^S K_s(x,c)\\
 &=(x-m_0)\sum_{s=0}^S\mathbf1_{R_s(c)}(x)\\
 &=(x-m_0)\mathbf1_{J_{S,j_S}}(x).
\end{aligned}
\]
This equality is exact at every half-open boundary.  A same-color alias is
zero already by Lemma~\ref{lem:mask-projection}.

Proposition~\ref{prop:refinement-independence} and the tower rule now yield
\[
 \theta(c)=\int_{J_{S,j_S}}(x-m_0)D(dx).
\]
For fixed $c$, $m_0,J_{S,j_S}$ are deterministic, and $X-m_0$ is integrable.
Splitting its integral gives
\[
\begin{aligned}
 \mu(D)-m_0
 &=\int_{\mathbb R}(x-m_0)D(dx)\\
 &=\int_{J_{S,j_S}}(x-m_0)D(dx)
  +\int_{J_{S,j_S}^c}(x-m_0)D(dx),
\end{aligned}
\]
so the second integral is the sole residual.

If $D(J_{0,j_0})=1$, Lemma~\ref{lem:four-color} shows that for every
$s\geq1$ neither the target nor any retained alias ring can be active.  A
mismatched retained color or branch makes $W_i(c)=0$ directly; when they
match, $F_i(X_i)=0$ and hence $\Delta Y_i=0$.  Thus every higher correction
is seedwise zero.  Nesting gives $J_{0,j_0}\subseteq J_{S,j_S}$, so the
outer residual is also zero and the level-zero identity supplies
$\theta(c)=\mu(D)-m_0$.  If $X_i=m_0$, the target level-zero coordinate is
zero as well, so every retained $\Delta Y_i$ and every $W_i(c)$ vanishes for
all seeds.

Thus Lemmas~\ref{lem:dither-moments} and \ref{lem:mask-projection},
Proposition~\ref{prop:level-inversion}, and the preceding finite telescope
account for every dither, mask, level, color, branch, and type factor and
leave exactly the displayed outer residual, with both exact baselines intact.
\end{proof}

\subsection{Target and alias distance with the all-ring activation ledger}

\begin{lemma}[Target-child exclusion forces residual distance]
\label{lem:target-distance}
Under Assumption~\ref{assump:parameter-domain} and
Lemma~\ref{lem:dyadic-ancestor}, let $c,x\in\mathbb R$ and
$1\leq s\leq S$.  If
\[
 x\in\mathcal R_{s,j_s,b_s}
 =J_{s,j_s}\setminus J_{s-1,j_{s-1}},
\]
then
\[
 |x-m_0|>\frac{h_s}{2},\qquad h_s<2|x-m_0|.
\]
This holds for both branches, signs of the path index, and all half-open
child-padding endpoints.
\end{lemma}

\begin{proof}
Define
\[
 a_s=j_0-2^sj_s
 =j_0-2^s\left\lfloor\frac{j_0}{2^s}\right\rfloor
 \in\{0,\ldots,2^s-1\}.
\]
Then
\[
 m_0=(j_s+\theta_s)h_s,
 \qquad \theta_s=\frac{a_s+1/2}{2^s},
\]
and
\[
 b_s=j_{s-1}-2j_s
 =\left\lfloor\frac{a_s}{2^{s-1}}\right\rfloor.
\]
These Euclidean-remainder identities remain valid for negative $j_0$.

If $b_s=0$, then $0\leq a_s\leq2^{s-1}-1$ and
\[
 \frac1{2^{s+1}}\leq\theta_s
 \leq\frac12-\frac1{2^{s+1}}.
\]
Since
$J_{s-1,j_{s-1}}=[(j_s-\tfrac12)h_s,(j_s+1)h_s)$,
\[
\begin{aligned}
 m_0-\frac{h_s}{2}&=(j_s-\tfrac12+\theta_s)h_s
 >(j_s-\tfrac12)h_s,\\
 m_0+\frac{h_s}{2}&=(j_s+\tfrac12+\theta_s)h_s
 <(j_s+1)h_s.
\end{aligned}
\]
Thus the closed radius-$h_s/2$ interval lies strictly inside the selected
child padding.

If $b_s=1$, then $2^{s-1}\leq a_s\leq2^s-1$ and
\[
 \frac12+\frac1{2^{s+1}}\leq\theta_s
 \leq1-\frac1{2^{s+1}}.
\]
Now $J_{s-1,j_{s-1}}=[j_sh_s,(j_s+\tfrac32)h_s)$ and
\[
\begin{aligned}
 m_0-\frac{h_s}{2}&=(j_s-\tfrac12+\theta_s)h_s>j_sh_s,\\
 m_0+\frac{h_s}{2}&=(j_s+\tfrac12+\theta_s)h_s
 <(j_s+\tfrac32)h_s.
\end{aligned}
\]
The same strict containment holds.  In either branch, a point in the target
ring lies outside the selected child padding and therefore outside this closed
interval.  Hence $|x-m_0|>h_s/2$.  The strict containment also handles the
asymmetric included-left and excluded-right endpoints.
\end{proof}

\begin{lemma}[Same-color aliases are separated from the path center]
\label{lem:alias-distance}
Under Assumption~\ref{assump:parameter-domain} and Lemmas~
\ref{lem:dyadic-ancestor} and \ref{lem:four-color}, let
$c,x\in\mathbb R$, $1\leq s\leq S$, and
$q\in\mathbb Z\setminus\{0\}$.  If
$x\in\mathcal R_{s,j_s+4q,b_s}$, then
\[
 |x-m_0|>2h_s,
 \qquad h_s<\frac{|x-m_0|}{2}.
\]
The conclusion is uniform over the retained branch, all index signs, and all
half-open boundaries.
\end{lemma}

\begin{proof}
Lemma~\ref{lem:dyadic-ancestor} places $m_0$ strictly in the path cell:
\[
 j_sh_s<m_0<(j_s+1)h_s.
\]
Every branch ring lies in its padding, so
\[
 \mathcal R_{s,j_s+4q,b_s}
 \subseteq[(j_s+4q-1)h_s,(j_s+4q+2)h_s).
\]
If $q\geq1$, activation implies
$x\geq(j_s+4q-1)h_s\geq(j_s+3)h_s$; since
$m_0<(j_s+1)h_s$, this gives $x-m_0>2h_s$.  Strictness remains at the
included left endpoint when $q=1$.

If $q\leq-1$, activation and the excluded right endpoint give
$x<(j_s+4q+2)h_s\leq(j_s-2)h_s$.  Since $m_0>j_sh_s$,
$m_0-x>2h_s$.  This is strict for $q=-1$ because the alias padding is
right-open and $m_0$ is strictly inside the path cell.  The two cases exhaust
all nonzero integers and are insensitive to the sign of $j_s$ or the branch.
\end{proof}

\begin{lemma}[Uniform all-ring activation ledger]
\label{lem:activation-ledger}
Under Assumptions~\ref{assump:parameter-domain} and
\ref{assump:precommitted-protocol}, Lemmas~\ref{lem:dyadic-ancestor},
\ref{lem:four-color}, \ref{lem:target-distance}, and
\ref{lem:alias-distance}, define
\[
 \Gamma_s(c)=\bigcup_{q\in\mathbb Z}
 \mathcal R_{s,j_s+4q,b_s},\qquad 1\leq s\leq S.
\]
For every $c,x$, with $r=|x-m_0|$,
\[
 x\in\Gamma_s(c)\quad\Longrightarrow\quad h_s\leq2r,
\]
and
\[
 \sum_{s=1}^S h_s^k\mathbf1\{x\in\Gamma_s(c)\}
 \leq C_{{\rm act},k}r^k,
 \qquad C_{{\rm act},k}=\frac{2^k}{1-2^{-k}}.
\]
At each level at most one ring is active, so the sum charges every retained
target and alias.  It is exactly zero when $r=0$ and for every
$x\in J_{0,j_0}$.
\end{lemma}

\begin{proof}
Every integer congruent to $j_s$ modulo four is uniquely $j_s+4q$.
Lemma~\ref{lem:four-color} makes the active index unique at each level.  For
$q=0$, Lemma~\ref{lem:target-distance} gives $h_s<2r$; for $q\neq0$,
Lemma~\ref{lem:alias-distance} gives the stronger $h_s<r/2$.  Thus no
retained ring is omitted and
\[
 \sum_{s=1}^S h_s^k\mathbf1\{x\in\Gamma_s(c)\}
 \leq\sum_{\substack{1\leq s\leq S\\h_s\leq2r}}h_s^k.
\]
If the index set is empty, the result is immediate.  Otherwise let $t$ be
its largest element.  Since $h_s=2^{s-t}h_t$,
\[
\begin{aligned}
 \sum_{\substack{1\leq s\leq S\\h_s\leq2r}}h_s^k
 &=\sum_{s=1}^t h_s^k
 =h_t^k\sum_{u=0}^{t-1}2^{-ku}\\
 &\leq\frac{h_t^k}{1-2^{-k}}
 \leq\frac{(2r)^k}{1-2^{-k}}
 =C_{{\rm act},k}r^k.
\end{aligned}
\]
The denominator is positive because $k>1$, and the constant is independent
of the path, branch, level count, decoder value, and sample point.

If $r=0$, no positive $h_s$ can satisfy $h_s\leq2r$, so the ledger is
zero.  If $x\in J_{0,j_0}$, nesting puts it in every selected child padding,
so the target ring is inactive; it also lies in $J_{s,j_s}$, which is
disjoint from every $q\neq0$ alias padding.  Hence the entire higher-level
ledger vanishes pointwise, including at all half-open endpoints.

The target and alias distance lemmas exhaust the retained ring indices, and
the geometric summation above therefore charges all higher-level target and
alias activations with the explicit constant $C_{{\rm act},k}$.
\end{proof}

\subsection{All-alias square ledgers and conditional variance}

For this subsection only, write
\[
 E_{0,q}(c)=\mathcal R_{0,j_0+4q},\qquad
 \Gamma_0(c)=\bigcup_{q\in\mathbb Z}E_{0,q}(c),
\]
and, for $1\leq s\leq S$,
\[
 E_{s,q}(c)=\mathcal R_{s,j_s+4q,b_s},\qquad
 \Gamma_s(c)=\bigcup_{q\in\mathbb Z}E_{s,q}(c).
\]
At fixed $(s,x,c)$, Lemma~\ref{lem:four-color} makes at most one term over
$q$ nonzero.

\begin{lemma}[Exact level-zero all-alias square ledger]
\label{lem:level-zero-square}
Under Assumptions~\ref{assump:independent-samples} and
\ref{assump:precommitted-protocol}, Proposition~
\ref{prop:refinement-independence}, and Lemmas~\ref{lem:four-color},
\ref{lem:query-measurability}, and \ref{lem:dither-moments}, every
$c,x\in\mathbb R$ satisfies
\[
\begin{aligned}
 &\mathbb E[W_i(c)^2\mathbf1\{L_i=0\}\mid X_i=x,c]\\
 &\quad=\frac{16h_0}{p_0}\sum_{q\in\mathbb Z}
 |x-m_{0,j_0+4q}|\mathbf1_{E_{0,q}(c)}(x)\\
 &\quad\leq\frac{24h_0^2}{p_0}\mathbf1_{\Gamma_0(c)}(x).
\end{aligned}
\]
The equality includes the target $q=0$ and every level-zero alias.
\end{lemma}

\begin{proof}
On $\{L_i=0,C_i=\kappa_0\}$, unique activity gives
\[
 F_i(x)=\sum_{q\in\mathbb Z}\rho_{i,0,j_0+4q}
 \frac{x-m_{0,j_0+4q}}{2h_0}\mathbf1_{E_{0,q}(c)}(x).
\]
Although $W_i(c)$ uses the target mask, its square contains
$\rho_{i,0,j_0}^2=1$.  Conditional on all masks, the dither-square identity
and unique activity yield
\[
 \mathbb E_{U_i}[\Delta Y_i^2]
 =\sum_{q\in\mathbb Z}
 \frac{|x-m_{0,j_0+4q}|}{4h_0}\mathbf1_{E_{0,q}(c)}(x).
\]
No mask cancellation occurs in this nonnegative square; an alias produces
the same absolute-amplitude term as the target.

Proposition~\ref{prop:refinement-independence} retains
$\Pr\{L_i=0\}=p_0$ and $\Pr\{C_i=\kappa_0\}=1/4$.  Therefore
\[
\begin{aligned}
 &p_0\cdot\frac14\left(\frac{16h_0}{p_0}\right)^2
 \sum_q\frac{|x-m_{0,j_0+4q}|}{4h_0}
  \mathbf1_{E_{0,q}(c)}(x)\\
 &\qquad=\frac{16h_0}{p_0}\sum_q|x-m_{0,j_0+4q}|
  \mathbf1_{E_{0,q}(c)}(x).
\end{aligned}
\]
The factor $p_0^{-1}$ remains explicit.  On an active level-zero ring,
$|x-m_{0,j}|\leq3h_0/2$ by Lemma~\ref{lem:four-color}; at most one $q$
is active.  This proves the factor-$24$ bound.  No branch or type probability
appears because neither variable exists at level zero.
\end{proof}

\begin{lemma}[Exact higher-level coordinate-mass all-alias square ledger]
\label{lem:higher-square}
Under Assumptions~\ref{assump:independent-samples} and
\ref{assump:precommitted-protocol}, Proposition~
\ref{prop:refinement-independence}, and Lemmas~\ref{lem:dyadic-ancestor},
\ref{lem:four-color}, \ref{lem:query-measurability}, and
\ref{lem:dither-moments}, every $c,x\in\mathbb R$ and
$1\leq s\leq S$ satisfy
\[
\begin{aligned}
 &\mathbb E[W_i(c)^2\mathbf1\{L_i=s\}\mid X_i=x,c]\\
 &\quad=\frac{64h_s}{p_s}\sum_{q\in\mathbb Z}
 |x-m_{s,j_s+4q}|\mathbf1_{E_{s,q}(c)}(x)
 +\frac{32d_s^2}{p_s}\mathbf1_{\Gamma_s(c)}(x)\\
 &\quad\leq\frac{128h_s^2}{p_s}\mathbf1_{\Gamma_s(c)}(x).
\end{aligned}
\]
Every target and alias appears in the exact sum; zero alias mean does not
remove its square.
\end{lemma}

\begin{proof}
On the retained color and branch, the two types are mutually exclusive, so
\[
 \left[4h_s\mathbf1\{T_i={\rm coord}\}
 +2d_s\mathbf1\{T_i={\rm mass}\}\right]^2
 =16h_s^2\mathbf1\{T_i={\rm coord}\}
 +4d_s^2\mathbf1\{T_i={\rm mass}\}.
\]
There is no coordinate-mass cross term.  For coordinate type, unique activity
and Lemma~\ref{lem:dither-moments} give
\[
 \mathbb E_{U_i}[\Delta Y_i^2]
 =\sum_q\frac{|x-m_{s,j_s+4q}|}{4h_s}
  \mathbf1_{E_{s,q}(c)}(x),
\]
whereas for mass type they give
\[
 \mathbb E_{U_i}[\Delta Y_i^2]
 =\frac12\mathbf1_{\Gamma_s(c)}(x).
\]
These quantities are independent of mask signs, and the target mask in
$W_i(c)^2$ squares to one even when an alias is active.

The exact probabilities of the level, retained color, retained branch, and
either type are $p_s,1/4,1/2,1/2$.  Hence the coordinate square is
\[
\begin{aligned}
 &p_s\cdot\frac14\cdot\frac12\cdot\frac12
 \left(\frac{16}{p_s}4h_s\right)^2\\[-0.2em]
 &\quad{}\times
 \sum_q\frac{|x-m_{s,j_s+4q}|}{4h_s}
  \mathbf1_{E_{s,q}(c)}(x)\\
 &=\frac{64h_s}{p_s}\sum_q|x-m_{s,j_s+4q}|\\[-0.2em]
 &\quad{}\times\mathbf1_{E_{s,q}(c)}(x),
\end{aligned}
\]
and the mass square is
\[
\begin{aligned}
 &p_s\cdot\frac14\cdot\frac12\cdot\frac12
 \left(\frac{16}{p_s}2d_s\right)^2
 \frac12\mathbf1_{\Gamma_s(c)}(x)\\
 &\qquad=\frac{32d_s^2}{p_s}\mathbf1_{\Gamma_s(c)}(x).
\end{aligned}
\]
Exactly one factor $p_s^{-1}$ remains after selection averaging.

On an active coordinate ring,
$|x-m_{s,j}|\leq3h_s/2$, so the coordinate term is at most
$96h_s^2p_s^{-1}\mathbf1_{\Gamma_s(c)}(x)$.  The displacement bound
$|d_s|\leq h_s$ makes the mass term at most
$32h_s^2p_s^{-1}\mathbf1_{\Gamma_s(c)}(x)$.  Their sum proves the stated
constant.  The exact mass term vanishes when $d_s=0$.
\end{proof}

\begin{proposition}[Pointwise all-level second-moment budget]
\label{prop:pointwise-square}
Under Assumptions~\ref{assump:independent-samples} and
\ref{assump:precommitted-protocol}, Lemmas~\ref{lem:level-zero-square},
\ref{lem:higher-square}, and \ref{lem:activation-ledger}, every fixed
$c,x\in\mathbb R$ satisfies
\[
 \mathbb E[W_i(c)^2\mid X_i=x,c]
 \leq Z_S\left[24h_0^k+128C_{{\rm act},k}|x-m_0|^k\right].
\]
This inequality sums all target and alias squares before taking expectation
over $X_i$.
\end{proposition}

\begin{proof}
The level events are mutually exclusive and exhaustive.  Lemmas~
\ref{lem:level-zero-square} and \ref{lem:higher-square} give
\[
\begin{aligned}
 \mathbb E[W_i(c)^2\mid X_i=x,c]
 &\leq\frac{24h_0^2}{p_0}\mathbf1_{\Gamma_0(c)}(x)
 +128\sum_{s=1}^S\frac{h_s^2}{p_s}
   \mathbf1_{\Gamma_s(c)}(x).
\end{aligned}
\]
For every level,
\[
 p_s^{-1}=Z_Sh_s^{k-2},
 \qquad \frac{h_s^2}{p_s}=Z_Sh_s^k.
\]
Substitution without asymptotic notation yields
\[
\begin{aligned}
 \mathbb E[W_i(c)^2\mid X_i=x,c]
 &\leq Z_S\left[24h_0^k\mathbf1_{\Gamma_0(c)}(x)
 +128\sum_{s=1}^S h_s^k\mathbf1_{\Gamma_s(c)}(x)\right]\\
 &\leq Z_S\left[24h_0^k
 +128C_{{\rm act},k}|x-m_0|^k\right].
\end{aligned}
\]
The last line uses only $\mathbf1_{\Gamma_0(c)}\leq1$ and
Lemma~\ref{lem:activation-ledger}.  Every within-level ring sum has at most
one active term and the level sum is finite, so no expectation or supremum is
interchanged with a countable ring sum.
\end{proof}

\begin{proposition}[Conditional variance and the single middle-regime factor]
\label{prop:conditional-variance}
Under Assumption~\ref{assump:moment-class}, Propositions~
\ref{prop:refinement-independence} and \ref{prop:pointwise-square}, and
Lemma~\ref{lem:recentered-moment}, suppose $a_k\geq200$ and fix a realized
localization output with $|c-\mu(D)|\leq50\sigma$.  Then
\[
 \operatorname{Var}(W_i(c)\mid c)
 \leq C_k^{\rm var}\sigma^kZ_S,
\]
where
\[
 C_k^{\rm var}=24a_k^k+128C_{{\rm act},k}C_k^{\rm rec}<\infty.
\]
The constant depends only on fixed $k$.  At $k=2$ the derivation contains
exactly one factor $Z_S=S+1$ and no second factor proportional to $S$.
\end{proposition}

\begin{proof}
Conditional on $c$, Proposition~\ref{prop:refinement-independence} leaves
$X_i$ with law $D$.  Integrating Proposition~\ref{prop:pointwise-square}
only after its pointwise level summation gives
\[
\begin{aligned}
 \mathbb E[W_i(c)^2\mid c]
 &\leq Z_S\left[24h_0^k+128C_{{\rm act},k}
   \int|x-m_0(c)|^kD(dx)\right]\\
 &=Z_S\left[24h_0^k+128C_{{\rm act},k}M_k(c)\right].
\end{aligned}
\]
On $\mathcal E_{\rm loc}$, Lemma~\ref{lem:recentered-moment} and
$h_0=a_k\sigma$ imply
\[
\begin{aligned}
 \mathbb E[W_i(c)^2\mid c]
 &\leq Z_S\left[24a_k^k\sigma^k
 +128C_{{\rm act},k}C_k^{\rm rec}\sigma^k\right]\\
 &=C_k^{\rm var}\sigma^kZ_S.
\end{aligned}
\]
This proves conditional square integrability, and
$\operatorname{Var}(W_i(c)\mid c)\leq\mathbb E[W_i(c)^2\mid c]$ gives the
variance bound.

At $k=2$,
\[
 Z_S=\sum_{s=0}^S h_s^0=S+1,
 \qquad p_s=\frac1{S+1},
 \qquad \frac{h_s^2}{p_s}=(S+1)h_s^2.
\]
Consequently Lemma~\ref{lem:higher-square} and the activation ledger, still
pointwise in $x$, give
\[
\begin{aligned}
 \sum_{s=1}^S\mathbb E[W_i(c)^2\mathbf1\{L_i=s\}\mid X_i=x,c]
 &\leq128(S+1)\sum_{s=1}^S h_s^2
   \mathbf1_{\Gamma_s(c)}(x)\\
 &\leq128C_{{\rm act},2}(S+1)|x-m_0|^2.
\end{aligned}
\]
The separate level-zero term is at most
\[
 24(S+1)h_0^2\mathbf1_{\Gamma_0(c)}(x).
\]
After expectation and recentering,
\[
\begin{aligned}
 \operatorname{Var}(W_i(c)\mid c)
 &\leq(S+1)\left[24h_0^2+128C_{{\rm act},2}M_2(c)\right]\\
 &\leq C_2^{\rm var}\sigma^2(S+1)
 =C_2^{\rm var}\sigma^2Z_S.
\end{aligned}
\]
Thus only the factor arising from $p_s^{-1}$ remains; the pointwise ledger
is uniform in $S$ and contributes no second factor.

If $D(J_{0,j_0})=1$, Lemma~\ref{lem:activation-ledger} makes every higher
target and alias square exactly zero, leaving only
Lemma~\ref{lem:level-zero-square}.  If $D\{m_0\}=1$, that exact formula
also vanishes and the conditional variance is zero.  These are exact
reductions, not remainder bounds.

The separate level-zero and higher-level square identities, their pointwise
sum, and the final conditional integration account for every alias square
and every importance factor.  In particular, the displayed middle-regime
calculation proves that the sole level cost is $Z_S=S+1$.
\end{proof}

\subsection{Outer scale, tail control, and three-regime normalization}

\begin{lemma}[Admissible dyadic outer scale and exact rounding]
\label{lem:outer-scale}
Under Assumption~\ref{assump:parameter-domain} and
Lemma~\ref{lem:recentered-moment}, take $a_k\geq200$ and define
\[
 b_k=\max\left\{a_k,(4C_k^{\rm rec})^{1/(k-1)}\right\}.
\]
Choose
\[
 0<c_k\leq\min\left\{\frac12,
   \left(\frac{b_k}{2a_k}\right)^{k-1}\right\}.
\]
For every $0<\epsilon\leq c_k\sigma$,
\[
 \frac{H_\star}{h_0}\geq2,qquad S\geq1,qquad
 H_\star\leq H<2H_\star.
\]
At $H_\star/h_0=2$ one has $S=1$ and $H=H_\star$; conversely, under
this admissibility condition $S=1$ only at $H_\star/h_0=2$.
\end{lemma}

\begin{proof}
The choices depend only on fixed $k$ and give
\[
 b_k\geq a_k,qquad b_k^{k-1}\geq4C_k^{\rm rec}.
\]
Put
\[
 x=\frac\sigma\epsilon,qquad
 R=\frac{H_\star}{h_0}=\frac{b_k}{a_k}x^{1/(k-1)}.
\]
The parameter domain and the choice of $c_k$ imply
\[
 R\geq\frac{b_k}{a_k}c_k^{-1/(k-1)}
 \geq\frac{b_k}{a_k}\frac{2a_k}{b_k}=2.
\]
Also $c_k\leq1/2$ gives $x\geq2$.  Let $q=\log_2R$.  The exact ceiling
$S=\lceil q\rceil$ gives $q\leq S<q+1$, whence
\[
 H_\star=h_0R\leq h_02^S=H<2h_0R=2H_\star.
\]
Since $R\geq2$, $S\geq1$.  If $R=2$, then $q=S=1$ and
$H=H_\star$.  Conversely $S=1$ gives $q\leq1$, while $R\geq2$ gives
$q\geq1$, so $R=2$.  No strict inequality in the parameter domain was used,
and the calculation includes $\epsilon=c_k\sigma$.
\end{proof}

\begin{proposition}[Outer residual at the target scale]\label{prop:tail-bound}
Under Assumptions~\ref{assump:parameter-domain} and
\ref{assump:moment-class}, Lemmas~\ref{lem:recentered-moment},
\ref{lem:dyadic-ancestor}, and \ref{lem:outer-scale}, and
Proposition~\ref{prop:mean-telescope}, every localization transcript in
$\mathcal E_{\rm loc}$ satisfies
\[
 |\mu(D)-m_0-\theta(c)|
 \leq\frac{M_k(c)}{H^{k-1}}
 \leq C_k^{\rm rec}\frac{\sigma^k}{H^{k-1}}
 \leq\frac\epsilon4.
\]
If $D(J_{S,j_S})=1$, the left side is exactly zero.  This holds in
particular when $D(J_{0,j_0})=1$, preserving the exact level-zero reduction.
\end{proposition}

\begin{proof}
For a fixed transcript in $\mathcal E_{\rm loc}$,
Proposition~\ref{prop:mean-telescope} gives
\[
 \mu(D)-m_0-\theta(c)=\int_{J_{S,j_S}^c}(x-m_0)D(dx).
\]
Writing $r=|x-m_0|$, Lemma~\ref{lem:dyadic-ancestor} gives the strict set
inclusion
\[
 \{x:x\notin J_{S,j_S}\}\subseteq\{x:r>H\},
\]
because every point with $r\leq H$, including both equality points, lies in
the padding.  Hence
\[
\begin{aligned}
 |\mu(D)-m_0-\theta(c)|
 &\leq\int r\mathbf1\{x\notin J_{S,j_S}\}D(dx)\\
 &\leq\int r\mathbf1\{r>H\}D(dx)\\
 &\leq\int\frac{r^k}{H^{k-1}}D(dx)
 =\frac{M_k(c)}{H^{k-1}}.
\end{aligned}
\]
The third line uses $r\leq r^k/H^{k-1}$ on $r>H$.  Lemma~
\ref{lem:recentered-moment} gives the second inequality.  Since
$H\geq H_\star$ and $k-1>0$,
\[
\begin{aligned}
 C_k^{\rm rec}\frac{\sigma^k}{H^{k-1}}
 &\leq C_k^{\rm rec}\frac{\sigma^k}{H_\star^{k-1}},\\
 H_\star^{k-1}
 &=b_k^{k-1}\sigma^{k-1}\frac\sigma\epsilon
 =b_k^{k-1}\frac{\sigma^k}{\epsilon}.
\end{aligned}
\]
Thus
\[
 C_k^{\rm rec}\frac{\sigma^k}{H^{k-1}}
 \leq C_k^{\rm rec}b_k^{1-k}\epsilon\leq\frac\epsilon4,
\]
where the last threshold is exactly $b_k^{k-1}\geq4C_k^{\rm rec}$.
No tail event is introduced.

The half-open boundary cases are also exact.  Since
$J_{S,j_S}=[(j_S-1)H,(j_S+2)H)$ and
$m_0=(j_S+\vartheta_S)H$ with $0<\vartheta_S<1$, an atom at the included
left boundary is more than $H$ from $m_0$ but remains in the truncated
integral; an atom at the excluded right boundary is more than $H$ away and
is charged by the tail inequality.  Atoms at $m_0\pm H$ lie strictly inside
the padding.  If $D(J_{S,j_S})=1$, the exact residual integral is zero before
any inequality.  Nesting gives this from $D(J_{0,j_0})=1$, and
Proposition~\ref{prop:mean-telescope} then leaves only the exact level-zero
mean.
\end{proof}

\begin{lemma}[Exact finite normalizer in the three moment regimes]
\label{lem:normalizer}
Under Assumption~\ref{assump:parameter-domain} and
Lemma~\ref{lem:outer-scale}, write $x=\sigma/\epsilon$.  For $k>2$,
\[
 Z_S=h_0^{2-k}\frac{1-2^{(2-k)(S+1)}}{1-2^{2-k}},
\]
and
\[
 a_k^{2-k}\sigma^{2-k}\leq Z_S
 \leq\frac{a_k^{2-k}}{1-2^{2-k}}\sigma^{2-k}.
\]
For $k=2$,
\[
 Z_S=S+1,
 \qquad
 \frac{\log x}{\log2}\leq Z_S
 <\frac{3+\log_2(b_2/a_2)}{\log2}\log x.
\]
For $1<k<2$,
\[
 Z_S=h_0^{2-k}\frac{2^{(2-k)(S+1)}-1}{2^{2-k}-1},
\]
\[
 H^{2-k}\leq Z_S\leq\frac{H^{2-k}}{1-2^{k-2}},
\]
and
\[
\begin{aligned}
 b_k^{2-k}\sigma^{2-k}x^{(2-k)/(k-1)}
 &\leq Z_S\\
 &<\frac{2^{2-k}b_k^{2-k}}{1-2^{k-2}}
 \sigma^{2-k}x^{(2-k)/(k-1)}.
\end{aligned}
\]
Every formula includes the minimum legal case $S=1$.
\end{lemma}

\begin{proof}
Since $h_s=2^sh_0$,
\[
 Z_S=h_0^{2-k}\sum_{s=0}^S2^{(2-k)s}.
\]
For $k>2$, let $\rho=2^{2-k}\in(0,1)$.  Finite multiplication gives
$(1-\rho)\sum_{s=0}^S\rho^s=1-\rho^{S+1}$, proving the exact formula.
Moreover,
\[
 1\leq\sum_{s=0}^S\rho^s
 \leq\sum_{s=0}^\infty\rho^s=\frac1{1-\rho}.
\]
Multiplication by
$h_0^{2-k}=a_k^{2-k}\sigma^{2-k}$ proves both bounds.  The denominator
is positive for fixed $k>2$; no uniform limit as $k\downarrow2$ is claimed.

At $k=2$, every term is one.  Put
\[
 R=\frac{H_\star}{h_0}=\frac{b_2}{a_2}x,
 \qquad q=\log_2R.
\]
The scale lemma gives $b_2/a_2\geq1$, $x\geq2$, and
$q\leq S<q+1$.  Therefore
\[
 Z_S=S+1\geq q+1\geq\log_2x=\frac{\log x}{\log2}.
\]
For the upper bound,
\[
\begin{aligned}
 Z_S&<q+2
 =\log_2x+\log_2(b_2/a_2)+2\\
 &\leq[3+\log_2(b_2/a_2)]\log_2x
 =\frac{3+\log_2(b_2/a_2)}{\log2}\log x.
\end{aligned}
\]
The absorption is explicit because
\[
 [3+\log_2(b_2/a_2)]\log_2x
 -[\log_2x+\log_2(b_2/a_2)+2]
 =[2+\log_2(b_2/a_2)](\log_2x-1)\geq0.
\]

For $1<k<2$, let $\rho=2^{2-k}>1$.  The finite identity
$\sum_{s=0}^S\rho^s=(\rho^{S+1}-1)/(\rho-1)$ gives the exact formula.
Factoring out the last term yields
\[
\begin{aligned}
 Z_S&=H^{2-k}\sum_{t=0}^S2^{(k-2)t},\\
 1&\leq\sum_{t=0}^S2^{(k-2)t}
 \leq\sum_{t=0}^\infty2^{(k-2)t}
 =\frac1{1-2^{k-2}}.
\end{aligned}
\]
Because $2-k>0$ and $H_\star\leq H<2H_\star$,
\[
 H_\star^{2-k}\leq H^{2-k}<2^{2-k}H_\star^{2-k},
\]
while
\[
 H_\star^{2-k}=b_k^{2-k}\sigma^{2-k}
 x^{(2-k)/(k-1)}.
\]
Combining the displays proves the last bounds.  The denominator is positive
for fixed $1<k<2$.  When $H_\star/h_0=2$, the scale lemma gives
$S=1,H=H_\star$, and the formulas reduce to the two-term sum
$h_0^{2-k}+H^{2-k}$, or to $2$ at $k=2$, with no endpoint exception.
\end{proof}

\begin{proposition}[Three-regime refinement scale and nondegeneracy]
\label{prop:scale-rate}
Under Assumption~\ref{assump:parameter-domain} and Lemmas~
\ref{lem:outer-scale} and \ref{lem:normalizer}, define
\[
 A_k=\frac{\sigma^kZ_S}{\epsilon^2},
 \qquad x=\frac\sigma\epsilon.
\]
For $k>2$,
\[
 a_k^{2-k}x^2\leq A_k
 \leq\frac{a_k^{2-k}}{1-2^{2-k}}x^2.
\]
For $k=2$,
\[
 \frac{x^2\log x}{\log2}\leq A_2
 <\frac{3+\log_2(b_2/a_2)}{\log2}x^2\log x.
\]
For $1<k<2$,
\[
 b_k^{2-k}x^{k/(k-1)}\leq A_k
 <\frac{2^{2-k}b_k^{2-k}}{1-2^{k-2}}x^{k/(k-1)}.
\]
Consequently
\[
 A_k\asymp_k
 \begin{cases}
 \sigma^2/\epsilon^2,&k>2,\\
 (\sigma^2/\epsilon^2)\log(\sigma/\epsilon),&k=2,\\
 (\sigma/\epsilon)^{k/(k-1)},&1<k<2,
 \end{cases}
\]
and $A_k$ has a positive explicit $k$-only lower bound.
\end{proposition}

\begin{proof}
For $k>2$, multiplying the normalizer bounds by
$\sigma^k/\epsilon^2$ uses
\[
 \frac{\sigma^k}{\epsilon^2}\sigma^{2-k}
 =\frac{\sigma^2}{\epsilon^2}=x^2,
\]
and proves the first pair with every geometric factor retained.  At $k=2$,
$A_2=x^2Z_S$, so multiplying the two exact logarithmic bounds proves the
second pair; the additive ceiling term has already been controlled.

For $1<k<2$,
\[
\begin{aligned}
 \frac{\sigma^k}{\epsilon^2}\sigma^{2-k}
 x^{(2-k)/(k-1)}
 &=x^2x^{(2-k)/(k-1)}\\
 &=x^{2+(2-k)/(k-1)}=x^{k/(k-1)},
\end{aligned}
\]
which proves the last pair with its rounding factor and denominator intact.

Finally $x\geq2$ by Lemma~\ref{lem:outer-scale}.  Thus
\[
 A_k\geq\underline A_k:=
 \begin{cases}
 4a_k^{2-k},&k>2,\\
 4,&k=2,\\
 2^{k/(k-1)}b_k^{2-k},&1<k<2.
 \end{cases}
\]
Every branch is finite, positive, and $k$-only.  All substitutions remain
valid at $\epsilon=c_k\sigma$, $S=1$, and $\lambda=\sigma$.  Constants are
allowed to deteriorate as a fixed off-middle $k$ approaches two; no continuity
through the separately evaluated middle case is asserted.

Together with Proposition~\ref{prop:tail-bound}, the scale and normalizer
lemmas therefore provide the exact $\epsilon/4$ residual certificate, every
dyadic boundary case, all three finite-sum formulas, and the positive
$k$-only lower interface for $A_k$ used by the count ceilings.
\end{proof}

\subsection{Conditional iid refinement and median amplification}

\begin{proposition}[Full-transcript conditional iid refinement law]
\label{prop:conditional-iid}
Under Assumptions~\ref{assump:independent-samples} and
\ref{assump:precommitted-protocol}, Propositions~
\ref{prop:refinement-independence}, \ref{prop:mean-telescope}, and
\ref{prop:conditional-variance}, conditional on the full localization
sigma-field $\mathscr L_{\rm loc}$, the family
$(W_i(c))_{i\in I_{\rm ref}}$ is iid and
\[
 \mathbb E[W_i(c)\mid\mathscr L_{\rm loc}]=\theta(c)
 \quad\text{almost surely}.
\]
On $\mathcal E_{\rm loc}$, its common conditional variance satisfies
\[
 v(c):=\operatorname{Var}(W_i(c)\mid\mathscr L_{\rm loc})
 \leq C_k^{\rm var}\sigma^kZ_S
 \quad\text{almost surely}.
\]
Both moments average over every refinement sample and all refinement seeds;
no assertion conditional on frozen refinement coins is made.
\end{proposition}

\begin{proof}
Let $\Xi_i$ denote the complete refinement tuple, assigning fixed dummy
values to $T_i,B_i$ on the level-zero branch.  Proposition~
\ref{prop:refinement-independence} supplies a version of the conditional law
with
\[
 \mathcal L((\Xi_i)_{i\in I_{\rm ref}}\mid\mathscr L_{\rm loc})
 =Q^{\otimes N_{\rm ref}},
\]
where $Q$ is the common refinement sample-and-seed law and is independent of
the localization transcript.  The statistic is the common measurable map
\[
 W_i(c)=w(c,\Xi_i),
\]
where the transcript enters only through its measurable output $c$ and the
resulting decoder path.  For a transcript realization $\ell$,
\[
 \mathcal L((W_i(c))_{i\in I_{\rm ref}}\mid\mathscr L_{\rm loc})(\ell)
 =\bigotimes_{i\in I_{\rm ref}}
 \mathcal L_Q(w(c(\ell),\Xi_i)).
\]
All factors are identical, proving conditional iid under the full transcript.

The conditional mean is
\[
 \int w(c(\ell),\xi)Q(d\xi)=\theta(c(\ell))
\]
by Proposition~\ref{prop:mean-telescope}.  On
$\mathcal E_{\rm loc}$ the same kernel gives
\[
 v(c(\ell))=\int(w(c(\ell),\xi)-\theta(c(\ell)))^2Q(d\xi),
\]
which Proposition~\ref{prop:conditional-variance} bounds by
$C_k^{\rm var}\sigma^kZ_S$.  This equality comes from the explicit
transcript-independent kernel and the dependence of $w$ only through $c$;
it does not invoke the generally false claim that enlarging a conditioning
sigma-field preserves conditional variance.

If $v(c)=0$, the displayed integral shows
$W_i(c)=\theta(c)$ $Q$-almost surely.  Fixed disjoint groups use disjoint
tuple families and are therefore conditionally independent.
\end{proof}

\begin{lemma}[One-group accuracy at the stochastic target]
\label{lem:block-accuracy}
Under Assumption~\ref{assump:parameter-domain} and
Propositions~\ref{prop:conditional-iid} and \ref{prop:scale-rate}, choose
\[
 \beta_k\geq16C_k^{\rm var},
 \qquad B_{\rm ref}=\lceil\beta_kA_k\rceil.
\]
For every fixed group $G_g$, almost surely on $\mathcal E_{\rm loc}$,
\[
 \Pr\left\{\left|\overline W_g(c)-\theta(c)\right|>\frac\epsilon2
 \ \middle|\ \mathscr L_{\rm loc}\right\}\leq\frac14.
\]
This includes zero variance and every value of the block ceiling.
\end{lemma}

\begin{proof}
Proposition~\ref{prop:scale-rate} gives $A_k>0$, and the explicit variance
constant is positive.  Thus the choice is legal and
\[
 1\leq B_{\rm ref},
 \qquad B_{\rm ref}\geq\beta_kA_k
 =\beta_k\frac{\sigma^kZ_S}{\epsilon^2}.
\]
Conditional on a transcript in $\mathcal E_{\rm loc}$, Proposition~
\ref{prop:conditional-iid} and the exact group size give
\[
 \mathbb E[\overline W_g(c)\mid\mathscr L_{\rm loc}]=\theta(c),
 \qquad
 \operatorname{Var}(\overline W_g(c)\mid\mathscr L_{\rm loc})
 =\frac{v(c)}{B_{\rm ref}}.
\]
Conditional Chebyshev at threshold $\epsilon/2$ yields
\[
\begin{aligned}
 \Pr\left\{|\overline W_g(c)-\theta(c)|>\frac\epsilon2
 \ \middle|\ \mathscr L_{\rm loc}\right\}
 &\leq\frac{4v(c)}{B_{\rm ref}\epsilon^2}\\
 &\leq\frac{4C_k^{\rm var}\sigma^kZ_S}
 {B_{\rm ref}\epsilon^2}\\
 &=\frac{4C_k^{\rm var}A_k}{B_{\rm ref}}
 \leq\frac{4C_k^{\rm var}}{\beta_k}\leq\frac14.
\end{aligned}
\]
If $v(c)=0$, every group mean equals $\theta(c)$ conditionally almost surely,
and the failure probability is exactly zero.  No path or cell union bound is
taken.
\end{proof}

\begin{lemma}[Odd-median amplification with exact confidence constants]
\label{lem:median-amplification}
Under Assumption~\ref{assump:parameter-domain}, Proposition~
\ref{prop:conditional-iid}, and Lemma~\ref{lem:block-accuracy}, choose
\[
 \alpha_k=4,
 \qquad G_\delta=2\left\lceil4\log\frac8\delta\right\rceil+1.
\]
Define the median as the $(G_\delta+1)/2$-th order statistic, retaining
repeated values.  Almost surely on $\mathcal E_{\rm loc}$,
\[
\begin{aligned}
 &\Pr\left\{
 \left|\mathop{\rm median}_{1\leq g\leq G_\delta}\overline W_g(c)
 -\theta(c)\right|>\frac\epsilon2
 \ \middle|\ \mathscr L_{\rm loc}\right\}\\
 &\hspace{10em}\leq e^{-G_\delta/8}\leq\frac\delta2.
\end{aligned}
\]
\end{lemma}

\begin{proof}
Write $G=G_\delta$ and define
\[
 I_g=\mathbf1\left\{|\overline W_g(c)-\theta(c)|>\frac\epsilon2\right\},
 \qquad S_G=\sum_{g=1}^GI_g.
\]
The fixed disjoint groups are conditionally independent by Proposition~
\ref{prop:conditional-iid}; hence the $I_g$ are conditionally independent
Bernoulli variables.  Lemma~\ref{lem:block-accuracy} gives, on
$\mathcal E_{\rm loc}$,
\[
 \mathbb E[I_g\mid\mathscr L_{\rm loc}]\leq\frac14,
 \qquad
 \mathbb E[S_G\mid\mathscr L_{\rm loc}]\leq\frac G4.
\]

The chosen $G$ is odd.  More generally, write $r=(G+1)/2$ and order the
group means as $z_{(1)}\leq\cdots\leq z_{(G)}$.  If
$z_{(r)}>\theta(c)+\epsilon/2$, the entries $r,\ldots,G$ give at least $r$
bad groups; if $z_{(r)}<\theta(c)-\epsilon/2$, the entries $1,\ldots,r$
do so.  Therefore
\[
 \left\{|z_{(r)}-\theta(c)|>\frac\epsilon2\right\}
 \subseteq\left\{S_G\geq\frac{G+1}{2}\right\}
 \subseteq\left\{S_G\geq\frac G2\right\}.
\]
An exact tie at either threshold is successful because failure is strict;
repeated values do not affect the fixed odd rank.

On the failure event,
\[
 S_G-\mathbb E[S_G\mid\mathscr L_{\rm loc}]
 \geq\frac G2-\frac G4=\frac G4.
\]
Conditional Hoeffding for independent variables in $[0,1]$ gives
\[
\begin{aligned}
 \Pr\left\{|z_{(r)}-\theta(c)|>\frac\epsilon2
 \ \middle|\ \mathscr L_{\rm loc}\right\}
 &\leq\Pr\left\{S_G-\mathbb E[S_G\mid\mathscr L_{\rm loc}]
 \geq\frac G4\ \middle|\ \mathscr L_{\rm loc}\right\}\\
 &\leq\exp\left(-\frac{2(G/4)^2}{G}\right)=e^{-G/8}.
\end{aligned}
\]
With $L_\delta=\log(8/\delta)>0$,
\[
 G_\delta=2\lceil4L_\delta\rceil+1\geq8L_\delta+1,
\]
and hence
\[
 e^{-G_\delta/8}\leq e^{-1/8}e^{-L_\delta}
 =e^{-1/8}\frac\delta8<\frac\delta2.
\]
This includes all $\delta\in(0,1/2)$, including the limiting smallest legal
ceiling as $\delta\uparrow1/2$.  If $v(c)=0$, every $I_g$ is conditionally
zero and the median failure probability is exactly zero.
\end{proof}

\begin{proposition}[Conditional refinement accuracy at the common target]
\label{prop:conditional-accuracy}
Under Assumption~\ref{assump:parameter-domain}, Proposition~
\ref{prop:tail-bound}, and Lemma~\ref{lem:median-amplification}, define
\[
 \mathcal E_{\rm ref}=\left\{
 \left|\mathop{\rm median}_{1\leq g\leq G_\delta}\overline W_g(c)
 -\theta(c)\right|\leq\frac\epsilon2\right\}.
\]
Then, almost surely,
\[
 \mathbf1_{\mathcal E_{\rm loc}}
 \Pr(\mathcal E_{\rm ref}^c\mid\mathscr L_{\rm loc})
 \leq\frac\delta2\mathbf1_{\mathcal E_{\rm loc}}.
\]
On $\mathcal E_{\rm loc}\cap\mathcal E_{\rm ref}$,
\[
 |\widehat\mu-\mu(D)|\leq\frac{3\epsilon}{4}<\epsilon.
\]
Equivalently,
\[
 \mathbf1_{\mathcal E_{\rm loc}}
 \Pr\left\{|\widehat\mu-\mu(D)|>\frac{3\epsilon}{4}
 \ \middle|\ \mathscr L_{\rm loc}\right\}
 \leq\frac\delta2\mathbf1_{\mathcal E_{\rm loc}},
\]
and the same inequality holds with threshold $\epsilon$.
\end{proposition}

\begin{proof}
Lemma~\ref{lem:median-amplification} gives the first conditional inequality
on each successful localization transcript.  On the same transcript,
Proposition~\ref{prop:tail-bound} gives the deterministic bound
\[
 |\theta(c)-(\mu(D)-m_0)|\leq\frac\epsilon4.
\]
On $\mathcal E_{\rm ref}$, the exact estimator definition and the triangle
inequality give
\[
\begin{aligned}
 |\widehat\mu-\mu(D)|
 &=\left|\mathop{\rm median}_{g}\overline W_g(c)
 -(\mu(D)-m_0)\right|\\
 &\leq\left|\mathop{\rm median}_{g}\overline W_g(c)-\theta(c)\right|
 +|\theta(c)-(\mu(D)-m_0)|\\
 &\leq\frac\epsilon2+\frac\epsilon4
 =\frac{3\epsilon}{4}<\epsilon.
\end{aligned}
\]
The stochastic error and deterministic outer residual each appear exactly
once.  There is no additional localization term because localization only
provides the event on which the recentered variance and tail bounds hold.
Thus total-error failure on $\mathcal E_{\rm loc}$ is contained in
$\mathcal E_{\rm ref}^c$, proving both indicator-valued inequalities without
yet integrating out localization.

If $D(J_{0,j_0(c)})=1$, Proposition~\ref{prop:mean-telescope} makes every
higher statistic seedwise zero and sets $\theta(c)=\mu(D)-m_0$, while
Proposition~\ref{prop:tail-bound} makes the residual zero.  The median proof
then amplifies only the level-zero correction.  If $D\{m_0\}=1$, every
$W_i(c)$, group mean, and median is zero, so
$\widehat\mu=m_0=\mu(D)$ exactly.

Proposition~\ref{prop:conditional-iid}, Lemmas~\ref{lem:block-accuracy} and
\ref{lem:median-amplification}, and this same-target inequality give the
full-localization-transcript conditional interface: failure at most
$\delta/2$ and total error $3\epsilon/4<\epsilon$, including zero variance
and the smallest legal odd group count.
\end{proof}

\subsection{Query legality, unconditional accuracy, and the rate bridge}

\begin{proposition}[Precommitted Borel query certificate]
\label{prop:query-legality}
Under Assumptions~\ref{assump:parameter-domain} and
\ref{assump:precommitted-protocol}, Proposition~\ref{prop:localization} and
Lemma~\ref{lem:query-measurability}, every query set used by the two-block
protocol is Borel and fixed before the first response.  This includes every
countable mask realization.  The later decoder choice of
$c,j_0,(j_s,b_s),m_s,d_s,R_s(c)$ and the coefficients in $W_i(c)$ changes
no query and introduces no new query.
\end{proposition}

\begin{proof}
For localization, Proposition~\ref{prop:localization} converts every
deterministic source bit into the Borel inverse image
$\mathcal B_i(R_{\rm loc})$.  Its codebook, clipped bins, minimum-index tie
rule, sample count, and degenerate seed depend only on known
$(\lambda,\sigma,\delta)$ and are fixed before any response.  In the
zero-query branch $N_{\rm loc}=0$, so the query family is empty.

For refinement index $i$, the complete seed
\[
 (L_i,C_i,T_i,B_i,(\rho_{i,s,j})_{s,j},U_i),
\]
with inapplicable level-zero variables omitted or given dummy values, is
drawn before any message.  Lemma~\ref{lem:query-measurability} proves that
the countable sum $F_i$ is jointly Borel, stabilizes pointwise because at
most one same-color ring is active, and obeys $|F_i(x)|\leq1$.  Hence
\[
 A_i=\{x\in\mathbb R:F_i(x)\geq U_i\}
\]
is Borel for every realized seed and dither.  Its formula contains the
global grid and pre-drawn level, color, type, branch, masks, and dither, but
no $c$, localization bit, selected path, or earlier response.

Only after both complete transcripts have arrived does the decoder evaluate
the Borel map $c\mapsto(j_0,j_s,b_s,m_s,d_s,R_s(c))$ and use it to retain
already-received bits and form $W_i(c)$.  No sample is queried then and no
fixed set $\mathcal B_i$ or $A_i$ changes.  This covers localization failure
and atoms at every half-open boundary.
\end{proof}

\begin{proposition}[Exact one-bit communication and fixed horizon]
\label{prop:one-bit-accounting}
Under Assumption~\ref{assump:precommitted-protocol} and
Proposition~\ref{prop:query-legality}, the protocol has deterministic
non-stopping horizon
\[
 n=N_{\rm loc}+N_{\rm ref},
 \qquad N_{\rm ref}=G_\delta B_{\rm ref}.
\]
Each of these independent samples transmits exactly one bit.  The baseline
$Y_i^0=\mathbf1\{0\geq U_i\}$ is computed from stored public randomness
and transmits no bit; decoder path selection and weighting transmit and query
nothing further.
\end{proposition}

\begin{proof}
The localization horizon is selected from known parameters.  In the
nontrivial branch each $i\in I_{\rm loc}$ sends only
\[
 Y_i^{\rm loc}=\mathbf1\{X_i\in\mathcal B_i(R_{\rm loc})\}\in\{0,1\}.
\]
In the trivial branch $N_{\rm loc}=0$, so no unused sample remains.

The values $B_{\rm ref},G_\delta$, the set $I_{\rm ref}$, and its equal
group partition are deterministic functions of known parameters and fixed
$k$-only choices.  Each $i\in I_{\rm ref}$ sends only
\[
 Y_i=\mathbf1\{X_i\in A_i\}\in\{0,1\}.
\]
The decoder already stores $U_i$ and computes $Y_i^0$, $\Delta Y_i$, the
selected path, $W_i(c)$, group means, their median, and $\widehat\mu$
locally.  Public seeds are not sample messages.  Hence every used sample
sends one bit and the horizon cannot stop or change with data, messages,
seeds, or localization success.
\end{proof}

\begin{proposition}[Unconditional uniform PAC conversion]
\label{prop:unconditional-pac}
Under Assumptions~\ref{assump:parameter-domain}--
\ref{assump:precommitted-protocol}, Propositions~\ref{prop:localization} and
\ref{prop:conditional-accuracy} imply, for every
$D\in\mathcal D(k,\lambda,\sigma)$,
\[
 \Pr_D\{|\widehat\mu-\mu(D)|>\epsilon\}
 \leq\frac\delta4+\frac\delta2
 =\frac{3\delta}{4}\leq\delta.
\]
The probability is unconditional over both sample blocks and every protocol
seed, and the inequality is uniform over the unrestricted moment class.
\end{proposition}

\begin{proof}
Fix an admissible $D$ and set
$\mathcal A=\{|\widehat\mu-\mu(D)|>\epsilon\}$.  The event
$\mathcal E_{\rm loc}$ is $\mathscr L_{\rm loc}$-measurable by
Proposition~\ref{prop:refinement-independence}, and Proposition~
\ref{prop:conditional-accuracy} states almost surely that
\[
 \mathbf1_{\mathcal E_{\rm loc}}
 \Pr_D(\mathcal A\mid\mathscr L_{\rm loc})
 \leq\frac\delta2\mathbf1_{\mathcal E_{\rm loc}}.
\]
The tower property with this measurable multiplier gives
\[
\begin{aligned}
 \Pr_D(\mathcal A\cap\mathcal E_{\rm loc})
 &=\mathbb E_D[\mathbf1_{\mathcal E_{\rm loc}}\mathbf1_{\mathcal A}]\\
 &=\mathbb E_D\left[\mathbf1_{\mathcal E_{\rm loc}}
   \Pr_D(\mathcal A\mid\mathscr L_{\rm loc})\right]\\
 &\leq\frac\delta2\Pr_D(\mathcal E_{\rm loc})\leq\frac\delta2.
\end{aligned}
\]
Proposition~\ref{prop:localization} gives
$\Pr_D(\mathcal E_{\rm loc}^c)\leq\delta/4$.  Hence
\[
\begin{aligned}
 \Pr_D(\mathcal A)
 &\leq\Pr_D(\mathcal E_{\rm loc}^c)
 +\Pr_D(\mathcal A\cap\mathcal E_{\rm loc})\\
 &\leq\frac\delta4+\frac\delta2
 =\frac{3\delta}{4}\leq\delta.
\end{aligned}
\]
All remaining refinement samples and seeds are integrated inside the
conditional probability.  No union bound over cells or paths is used.  The
estimator is defined on $\mathcal E_{\rm loc}^c$, and that whole event is
paid.  Since $D$ was arbitrary and every constant is class-uniform, the
supremum over $\mathcal D(k,\lambda,\sigma)$ follows.
\end{proof}

\begin{proposition}[All-ceiling technical sample bound]
\label{prop:technical-count}
Under Assumption~\ref{assump:parameter-domain}, Propositions~
\ref{prop:localization} and \ref{prop:scale-rate}, and Lemmas~
\ref{lem:outer-scale}, \ref{lem:block-accuracy}, and
\ref{lem:median-amplification}, there is a finite $\widetilde C_k$ depending
only on fixed $k$ such that
\[
 N_{\rm loc}+N_{\rm ref}
 \leq\widetilde C_k\left[
 \log\frac\lambda\sigma+A_k\log\frac1\delta\right].
\]
The bound includes the dyadic, localization, block, and group ceilings and
absorbs $1+\log(4/\delta)$ through a positive $k$-only lower bound on $A_k$.
\end{proposition}

\begin{proof}
Use the complete legal design
\[
 a_k\geq200,
 \qquad b_k=\max\{a_k,(4C_k^{\rm rec})^{1/(k-1)}\},
\]
\[
 0<c_k\leq\min\left\{\frac12,
 \left(\frac{b_k}{2a_k}\right)^{k-1}\right\},
 \qquad \beta_k=16C_k^{\rm var},
 \qquad \alpha_k=4.
\]
Every quantity depends only on fixed $k$.  Lemma~\ref{lem:outer-scale}
proves from the exact ceiling that
\[
 \frac{H_\star}{h_0}\geq2,qquad S\geq1,qquad
 H_\star\leq H<2H_\star.
\]
The threshold $b_k^{k-1}\geq4C_k^{\rm rec}$ gives outer bias at most
$\epsilon/4$ by Proposition~\ref{prop:tail-bound}; the choices of
$\beta_k$ and $\alpha_k$ give group failure at most $1/4$ and conditional
median failure at most $\delta/2$.

Put
\[
 L_\delta=\log\frac1\delta,
 \qquad \ell_{\lambda,\sigma}=\log\frac\lambda\sigma.
\]
Assumption~\ref{assump:parameter-domain} gives
$L_\delta\geq\log2>0$ and $\ell_{\lambda,\sigma}\geq0$.
Proposition~\ref{prop:scale-rate} proves
\[
 A_k\geq\underline A_k>0,
 \qquad
 \underline A_k=
 \begin{cases}
 4a_k^{2-k},&k>2,\\
 4,&k=2,\\
 2^{k/(k-1)}b_k^{2-k},&1<k<2.
 \end{cases}
\]
Every branch is positive and depends only on fixed $k$.

The block ceiling obeys
\[
\begin{aligned}
 \beta_kA_k
 &\leq B_{\rm ref}=\lceil\beta_kA_k\rceil
 <\beta_kA_k+1\\
 &\leq(\beta_k+\underline A_k^{-1})A_k.
\end{aligned}
\]
Set $C_{B,k}=\beta_k+\underline A_k^{-1}$.  The odd-group ceiling obeys
\[
 G_\delta=2\left\lceil4\log\frac8\delta\right\rceil+1
 <8\log\frac8\delta+3.
\]
Because
\[
 \log\frac8\delta=L_\delta+3\log2\leq4L_\delta,
 \qquad 3\leq\frac3{\log2}L_\delta,
\]
we have
\[
 G_\delta<C_GL_\delta,
 \qquad C_G=32+\frac3{\log2}.
\]
Thus both refinement ceilings are retained in
\[
 N_{\rm ref}=G_\delta B_{\rm ref}
 \leq C_GC_{B,k}A_kL_\delta.
\]

Proposition~\ref{prop:localization}, which already includes the source bin
and code-length ceilings, gives
\[
 N_{\rm loc}\leq C_{{\rm loc},k}
 [1+\ell_{\lambda,\sigma}+\log(4/\delta)],
 \qquad C_{{\rm loc},k}=10001.
\]
Furthermore,
\[
 \log\frac4\delta=L_\delta+2\log2\leq3L_\delta,
 \qquad 1\leq\frac{L_\delta}{\log2}.
\]
With $C_L=3+1/\log2$,
\[
 1+\log\frac4\delta
 \leq C_LL_\delta
 \leq\frac{C_L}{\underline A_k}A_kL_\delta.
\]
This is the required localization-confidence absorption.  Therefore
\[
 N_{\rm loc}\leq C_{{\rm loc},k}\ell_{\lambda,\sigma}
 +\frac{C_{{\rm loc},k}C_L}{\underline A_k}A_kL_\delta.
\]
Adding both blocks yields
\[
\begin{aligned}
 N_{\rm loc}+N_{\rm ref}
 &\leq C_{{\rm loc},k}\ell_{\lambda,\sigma}\\
 &\quad+\left[
 \frac{C_{{\rm loc},k}C_L}{\underline A_k}+C_GC_{B,k}
 \right]A_kL_\delta.
\end{aligned}
\]
The proposition follows with
\[
 \widetilde C_k=\max\left\{C_{{\rm loc},k},
 \frac{C_{{\rm loc},k}C_L}{\underline A_k}+C_GC_{B,k}\right\}.
\]
No additive term remains outside
$\ell_{\lambda,\sigma}+A_kL_\delta$.
\end{proof}

\begin{proposition}[Rate Specialization Bridge]
\label{prop:rate-specialization-bridge}
Under Assumptions~\ref{assump:parameter-domain}--
\ref{assump:precommitted-protocol}, take
\[
 a_k\geq200,
 \quad b_k=\max\{a_k,(4C_k^{\rm rec})^{1/(k-1)}\},
 \quad \beta_k=16C_k^{\rm var},
 \quad \alpha_k=4,
\]
and choose
\[
 0<c_k\leq\min\left\{\frac12,
 \left(\frac{b_k}{2a_k}\right)^{k-1}\right\}.
\]
Then all scale, tail, block, group, confidence, and probability conditions
of Propositions~\ref{prop:tail-bound}, \ref{prop:conditional-accuracy},
\ref{prop:unconditional-pac}, and \ref{prop:technical-count} hold, and there
is a finite $C_k>0$, depending only on fixed $k$, such that
\[
 N_{\rm loc}+N_{\rm ref}
 \leq C_kr_k(\lambda,\sigma,\epsilon,\delta).
\]
The probability mode is unconditional over all samples and seeds, the horizon
is fixed and non-stopping, the norm is absolute value on $\mathbb R$, and the
hidden constant is independent of $D,\lambda,\sigma,\epsilon,\delta$ and all
realized protocol objects.
\end{proposition}

\begin{proof}
The displayed choices are $k$-only.  Lemma~\ref{lem:outer-scale} verifies
\[
 H_\star/h_0\geq2,qquad S\geq1,qquad H_\star\leq H<2H_\star,
\]
and the choice of $b_k$ verifies the tail threshold
\[
 C_k^{\rm rec}\frac{\sigma^k}{H^{k-1}}leq\frac\epsilon4.
\]
Lemmas~\ref{lem:block-accuracy} and \ref{lem:median-amplification} verify the
group threshold $\beta_k=16C_k^{\rm var}$, the odd-group choice
$\alpha_k=4$, conditional group failure at most $1/4$, and conditional
median failure at most $e^{-G_\delta/8}\leq\delta/2$.  Proposition~
\ref{prop:conditional-accuracy} then gives total conditional error
$\epsilon/2+\epsilon/4=3\epsilon/4<\epsilon$.

The block, group, localization, and confidence ceilings were retained in
Proposition~\ref{prop:technical-count}.  In particular,
\[
 A_k\geq\underline A_k>0,
 \quad B_{\rm ref}<(\beta_k+\underline A_k^{-1})A_k,
 \quad G_\delta<\left(32+\frac3{\log2}\right)\log\frac1\delta,
\]
and
\[
 1+\log\frac4\delta
 \leq\frac{3+1/\log2}{\underline A_k}
 A_k\log\frac1\delta.
\]
Consequently
\[
 N_{\rm loc}+N_{\rm ref}
 \leq\widetilde C_k\left[
 \log\frac\lambda\sigma+A_k\log\frac1\delta\right].
\]

Define
\[
 Q_k(\sigma,\epsilon)=
 \begin{cases}
 \sigma^2/\epsilon^2,&k>2,\\
 (\sigma^2/\epsilon^2)\log(\sigma/\epsilon),&k=2,\\
 (\sigma/\epsilon)^{k/(k-1)},&1<k<2.
 \end{cases}
\]
Proposition~\ref{prop:scale-rate} gives, without asymptotic substitution,
\[
 A_k\leq U_kQ_k(\sigma,\epsilon),
\]
where
\[
 U_k=
 \begin{cases}
 \dfrac{a_k^{2-k}}{1-2^{2-k}},&k>2,\\[0.8em]
 \dfrac{3+\log_2(b_2/a_2)}{\log2},&k=2,\\[0.8em]
 \dfrac{2^{2-k}b_k^{2-k}}{1-2^{k-2}},&1<k<2.
 \end{cases}
\]
For each fixed regime the denominator is positive and $U_k$ is finite and
$k$-only.  The middle expression comes from the exact identity $Z_S=S+1$
and contains exactly one factor $\log(\sigma/\epsilon)$.  Substitution gives
\[
\begin{aligned}
 N_{\rm loc}+N_{\rm ref}
 &\leq\widetilde C_k\left[
 \log\frac\lambda\sigma+U_kQ_k(\sigma,\epsilon)\log\frac1\delta\right]\\
 &\leq\widetilde C_k\max\{1,U_k\}
 \left[\log\frac\lambda\sigma
 +Q_k(\sigma,\epsilon)\log\frac1\delta\right]\\
 &=C_kr_k(\lambda,\sigma,\epsilon,\delta),
\end{aligned}
\]
with $C_k=\widetilde C_k\max\{1,U_k\}$.  Its dependence is only through
fixed-$k$ design and geometric constants.

The probability conversion is also explicit.  Proposition~
\ref{prop:unconditional-pac} integrates the complete localization sigma-field:
\[
 \Pr_D(\mathcal A\cap\mathcal E_{\rm loc})
 =\mathbb E_D[\mathbf1_{\mathcal E_{\rm loc}}
 \Pr_D(\mathcal A\mid\mathscr L_{\rm loc})]\leq\frac\delta2,
\]
and then
\[
 \Pr_D(\mathcal A)\leq\frac\delta4+\frac\delta2
 =\frac{3\delta}{4}\leq\delta.
\]
No path union bound appears.  The substitutions remain valid at
$\lambda=\sigma$, $\epsilon=c_k\sigma$, $S=1$, and as
$\delta\uparrow1/2$, because $\log(1/\delta)\geq\log2$ supplies the explicit
absorption.  No continuity through $k=2$ is asserted.
\end{proof}

\begin{proposition}[Exact supported-cell and point-mass baselines]
\label{prop:baseline-reductions}
Under Propositions~\ref{prop:ring-partition}, \ref{prop:mean-telescope},
\ref{prop:conditional-variance}, \ref{prop:tail-bound}, and
\ref{prop:conditional-accuracy}, and Lemmas~\ref{lem:four-color} and
\ref{lem:activation-ledger}, fix any decoder output $c$.  If
$D(J_{0,j_0(c)})=1$, every retained higher correction is seedwise zero, the
higher activation and variance charges vanish, the outer residual is zero,
and the estimator is exactly the level-zero unbiased dither correction with
the same fixed median aggregation.  If $D\{m_0(c)\}=1$, every $W_i(c)=0$
seedwise and $\widehat\mu=m_0(c)=\mu(D)$.
\end{proposition}

\begin{proof}
Suppose $D(J_{0,j_0(c)})=1$.  Path nesting places $J_{0,j_0}$ inside every
selected child padding, so every higher target ring is inactive.  Four-color
separation excludes every retained alias.  A mismatched decoder color or
branch indicator makes $W_i(c)=0$ directly.  When the indicators match, the
sample lies on no retained ring, so $F_i(X_i)=0$, $Y_i=Y_i^0$, and again
$W_i(c)=0$.  This is seedwise equality, not cancellation in expectation.

Thus the higher activation and square ledgers are identically zero.
Proposition~\ref{prop:mean-telescope} reduces to
$\theta(c)=\mu(D)-m_0$ using level zero alone.  Since
$J_{0,j_0}\subseteq J_{S,j_S}$, the outer residual is zero before its
moment bound.  Hence
\[
 \widehat\mu=m_0+\mathop{\rm median}_{1\leq g\leq G_\delta}
 \left[\frac1{B_{\rm ref}}\sum_{i\in G_g}
 W_i(c)\mathbf1\{L_i=0\}\right].
\]
This is exactly the same level-zero dither estimate and fixed median; the
conditional concentration proof adds no higher-ring or tail remainder.

If $D\{m_0(c)\}=1$, the retained target level-zero amplitude
$(X_i-m_0)/(2h_0)$ is zero, while every higher ring is inactive.  Whenever a
decoder coefficient is nonzero, $F_i(X_i)=0$ and $\Delta Y_i=0$.  Thus
$W_i(c)=0$ for every sample and every public seed.  All group means and the
median are zero, and the point-mass condition gives
$\widehat\mu=m_0(c)=\mu(D)$ with no artificial importance, tail,
localization, or confidence residual.
\end{proof}

\subsection{Proof of the Main Theorem}

\begin{proof}[Proof of Theorem~\ref{thm:main}]
Choose $a_k,b_k,c_k,\alpha_k,\beta_k$ as in Proposition~
\ref{prop:rate-specialization-bridge}.  Proposition~\ref{prop:localization}
provides the always-defined localization output and its fixed cost, while
Lemma~\ref{lem:query-measurability} and Proposition~
\ref{prop:query-legality} prove that all localization and refinement queries
are Borel and simultaneously precommitted.  Proposition~
\ref{prop:one-bit-accounting} proves exactly one transmitted bit per used
independent sample and the deterministic non-stopping horizon
$N_{\rm loc}+N_{\rm ref}$.

Proposition~\ref{prop:mean-telescope} identifies the exact truncated target,
Proposition~\ref{prop:conditional-variance} controls its refinement variance,
and Proposition~\ref{prop:tail-bound} bounds the sole outer residual by
$\epsilon/4$.  Proposition~\ref{prop:conditional-accuracy} gives conditional
error at most $3\epsilon/4<\epsilon$ with conditional failure at most
$\delta/2$.  Proposition~\ref{prop:unconditional-pac} integrates the full
localization sigma-field and combines this with localization failure
$\delta/4$, proving
\[
 \sup_{D\in\mathcal D(k,\lambda,\sigma)}
 \Pr_D\{|\widehat\mu-\mu(D)|>\epsilon\}\leq\delta.
\]

Proposition~\ref{prop:rate-specialization-bridge} checks every auxiliary
threshold and ceiling and proves
\[
 N_{\rm loc}+N_{\rm ref}
 \leq C_kr_k(\lambda,\sigma,\epsilon,\delta)
\]
with $C_k$ depending only on fixed $k$, including exactly one
$\log(\sigma/\epsilon)$ in the middle regime.  Finally,
Proposition~\ref{prop:baseline-reductions} gives the stated supported-cell
level-zero reduction and exact point-mass recovery.  These conclusions are
precisely all clauses of Theorem~\ref{thm:main}.
\end{proof}
